%% -*- coding: utf-8 -*-
\documentclass[12pt,pagesize,paper=192mm:108mm,landscape]{scrbook} 
%1920x1080 1280x720
\areaset[current]{192mm}{108mm}
\usepackage{calc}
\usepackage[T2A]{fontenc}
\usepackage[utf8]{inputenc}
\usepackage[english,russian]{babel}
\usepackage{microtype}
\usepackage{misccorr}
\usepackage{cmap}
%\usepackage[unicode=true]{hyperref}
\usepackage{graphicx}
\usepackage{amssymb}
\usepackage{amsmath}
%\usepackage{srcltx}
\usepackage{textcomp}
\usepackage{xspace}
%научные символы и смайлики \smiley \frownie
\usepackage{wasysym}
\usepackage{ccicons}
\begin{document}
\begin{titlepage}
  \vspace*{-0.5em}
  \begin{center}    
    \hspace*{3em}
    \begin{minipage}[t]{1.5em}
      \includegraphics[width=\textwidth]{../BINP-logo}
    \end{minipage}\hfill
    \begin{minipage}{0.115\linewidth}
    \includegraphics[width=\textwidth]{../NSU-logo}
    \end{minipage}
    \hfill
    \hspace*{4.5em}

    Кафедра теоретической физики физического факультета НГУ

    \Large
    Профессор Грабовский А.\,В.

    \huge
    \textbf{Стандартная модель}

    \Large
    Лекция № 1
    \vfill

    \small
    \begin{minipage}{0.9\linewidth}
      Группы $U(1)_{Y_W}$ слабого гиперзаряда и $SU(2)_{I_W}$ слабого
      изоспина. Поля $W_{1,2,3}$ и $B$, заряды $g1$ и $g2$, удлинение
      производных, тензоры напряженности полей, лагранжианы
      Янга"--~Миллса для этих полей. Гиперзаряд, заряд и изоспин для
      частиц в рамках одного поколения. Три поколения частиц. Действие
      удлиненной производной на фермионы. Лагранжиан взаимодействия
      калибровочных полей и фермионов. Лагранжиан и потенциал поля
      Хиггса. Гиперзаряд, заряд и изоспин компонент поля
      Хиггса. Действие удлиненной производной на поле
      Хиггса. Вакуумное среднее поля Хиггса. Поля $Z$"=бозона и
      $A$"=фотона. Выражение масс $W$, $Z$ и $A$ через заряды $g1$, $g2$ и
      вакуумное среднее поля Хиггса. Угол Вайнберга. Заряженный ток "---
      взаимодействие фермионов с $W^{\pm}$. Нейтральный ток "--- взаимодействие
      фермионов с фотонов и $Z$-бозоном. Константы $g_V$ и $g_A$ для различных
      фермионов.
    \end{minipage}
    \vfill

    \normalsize \ccbysa\hspace{0.5em}  Новосибирск 2024
  \end{center}
\end{titlepage}
\end{document}
